\chapter*{LAMPIRAN}
\addcontentsline{toc}{chapter}{LAMPIRAN}
\thispagestyle{mainbody}

\buatlampiran{Jadwal Penelitian}

\noindent
{
\setstretch{1.0}
\begin{tblr}{
    width = \linewidth,
    colspec = {c X[l] *{5}{c}},
    hlines, vlines,
    rows = {m},
    rowhead = 3,
    row{1,2,3} = {font=\bfseries, halign=c, valign=m},
}
    \SetCell[r=3]{c} No & \SetCell[r=3]{c} Jenis Kegiatan & \SetCell[c=5]{c} Waktu Pelaksanaan & & & & \\
    & & \SetCell[c=4]{c} 2025 & & & & 2026 \\
    & & Sept & Okt & Nov & Des & Jan \\
    1 & Identifikasi Masalah & \SetCell{bg=green!25} & \SetCell{bg=green!25} & & & \\
    2 & Studi Literatur & \SetCell{bg=green!25} & \SetCell{bg=green!25} & & & \\
    3 & Pengumpulan Data & \SetCell{bg=green!25} & \SetCell{bg=green!25} & & & \\
    4 & Perancangan Sistem & & & \SetCell{bg=green!25} & & \\
    5 & Pengembangan Sistem & & & \SetCell{bg=green!25} & \SetCell{bg=green!25} & \\
    6 & Pengujian dan Evaluasi & & & & \SetCell{bg=green!25} & \SetCell{bg=green!25} \\
    7 & Penyusunan Laporan & & & & \SetCell{bg=green!25} & \SetCell{bg=green!25} \\
\end{tblr}
}

\buatlampiran{Repositori Kode Sumber Hasil Penelitian}

\noindent Seluruh kode implementasi penelitian hingga hasil pengujian dapat diakses pada tautan dan pemindaian kode QR berikut: \url{https://github.com/andyathsid/thesis}

\vspace{0.5cm}

\noindent \includegraphics[width=4.5cm]{lampiran/repository_monorepo.png}  

\buatlampiran{Hasil Pengujian Satu Tahap pada Dataset \textit{In-Distribution} (ID)}

\noindent Seluruh rekam jejak dan hasil pengujian satu tahap pada dataset \textit{In-Distribution} (ID) melalui platform LangSmith dapat diakses menggunakan pemindaian kode QR berikut.

\vspace{0.5cm}

\noindent 
\includegraphics[width=\linewidth]{lampiran/pengujian-satu-tahap-id.png}

\buatlampiran{Hasil Pengujian Satu Tahap pada Dataset \textit{Out-of-Distribution} (OOD)}

\noindent Seluruh rekam jejak dan hasil pengujian satu tahap pada dataset \textit{Out-of-Distribution} (OOD) melalui platform LangSmith dapat diakses menggunakan pemindaian kode QR berikut.

\vspace{0.5cm}

\noindent 
\includegraphics[width=\linewidth]{lampiran/pengujian-satu-tahap-ood.png}

\buatlampiran{Hasil Pengujian Multi-Tahap}

\noindent Seluruh rekam jejak dan hasil pengujian multi-tahap melalui platform LangSmith dapat diakses menggunakan pemindaian kode QR berikut.

\vspace{0.5cm}

\noindent 
\includegraphics[width=4.5cm]{lampiran/pengujian-multi-tahap.png}

\buatlampiran{Hasil Uji Similaritas Turnitin}

\noindent Hasil pemeriksaan similaritas dokumen tugas akhir menggunakan layanan Turnitin dilampirkan sebagai berikut.

\vspace{0.5cm}

\noindent \includegraphics[width=0.8\linewidth]{lampiran/uji-similaritas-turnitin.png}

\buatlampiran{Hasil Uji Deteksi AI}

\noindent Hasil pemeriksaan deteksi teks kecerdasan buatan (AI) pada dokumen tugas akhir menggunakan layanan GPTZero dilampirkan sebagai berikut.

\vspace{0.5cm}

\noindent \includegraphics[width=0.8\linewidth]{lampiran/uji-deteksi-ai.jpg}